\worldchapter{\trans{foes_title}}{appendix-foes}
\label{ch:appendix-foes}

\begin{multicols}{2}


    \section{Undead}



    \subsection*{Vampire spawn}

    The vampire spawn is a lower vampire. Without much influence, this is a figure which has limited powers, and may be left to its own devices.

    \begin{tabular}{@{}ll@{}}
        \textbf{\trans{wounds_label}:} & 10 \\
        \textbf{\trans{movement_label}:} & 4 \\
        \textbf{\trans{strength_label}:} & 2 \\
        \textbf{\trans{dexterity_label}:} & 2 \\
        \textbf{\trans{mind_label}:} & 2 \\
        \textbf{\trans{resistances_label}:} & ['Cold', 'Water'] \\
    \end{tabular}

    \textbf{Bite} (8)\\
    Bleeding 1\\
    \textbf{Black Magic} (6)\\
    The Vampire Spawn casts a spell of black magic. The successes of this roll represent the \textbf{power} of the spell.\\


    \subsection*{Zombie}

    The zombie is a tragic figure that appears in many stories. A human being who is kept alive in a special and unnatural way after his death. Zombies have a brain the size of a pea, they know little more than the desire for blood. And brains. So, if pop culture is to be believed.

    \begin{tabular}{@{}ll@{}}
        \textbf{\trans{wounds_label}:} & 6 \\
        \textbf{\trans{movement_label}:} & 3 \\
        \textbf{\trans{strength_label}:} & 3 \\
        \textbf{\trans{dexterity_label}:} & 1 \\
        \textbf{\trans{mind_label}:} & 1 \\
        \textbf{\trans{resistances_label}:} & ['Fire', 'Poison'] \\
    \end{tabular}

    \textbf{Undead Grip} (4)\\
    Poison 1\\
    \textbf{Bite} (5)\\
    Poison 1\\


    \subsection*{Mummy}

    The living mummy is undead. It has been magically brought to life. Nothing of her spirit remains, all she seeks is to take the life of her victims. She is usually unarmed, but her curse poisons her victims.

    \begin{tabular}{@{}ll@{}}
        \textbf{\trans{wounds_label}:} & 8 \\
        \textbf{\trans{movement_label}:} & 3 \\
        \textbf{\trans{strength_label}:} & 3 \\
        \textbf{\trans{dexterity_label}:} & 1 \\
        \textbf{\trans{mind_label}:} & 2 \\
        \textbf{\trans{resistances_label}:} & ['Fire', 'Poison'] \\
    \end{tabular}

    \textbf{Grip} (6)\\
    Piercing 0\\
    \textbf{Curse} (4)\\
    Piercing 3, Poison 2\\


    \subsection*{Skelett}

    Ein wandelndes Skelett, von dunkler Magie belebt.

    \begin{tabular}{@{}ll@{}}
        \textbf{\trans{wounds_label}:} & 4 \\
        \textbf{\trans{movement_label}:} & 3 \\
        \textbf{\trans{strength_label}:} & 2 \\
        \textbf{\trans{dexterity_label}:} & 2 \\
        \textbf{\trans{mind_label}:} & 2 \\
    \end{tabular}

    \textbf{Knochengriff} (8)\\
    Vergiftet 1\\
    \textbf{Kalter Atem} (6)\\
    Geschockt 2, 5 Meter Reichweite\\


    \section{Animal}



    \subsection*{Vulture}

    On rising thermals, the vulture circles patiently, watching for signs of weakness below. It does not hunt in the traditional sense—waiting instead for the inevitable, drawn to death, decay, and the promise of easy sustenance. Its presence is a silent omen, often arriving before the end has fully come. On the ground, it moves with awkward persistence, tearing and feeding with grim efficiency. In darker interpretations, it feels like a herald of fate, appearing wherever life falters, as if summoned by the dying.

    \begin{tabular}{@{}ll@{}}
        \textbf{\trans{wounds_label}:} & 2 \\
        \textbf{\trans{movement_label}:} & 6 \\
        \textbf{\trans{strength_label}:} & 1 \\
        \textbf{\trans{dexterity_label}:} & 3 \\
        \textbf{\trans{mind_label}:} & 2 \\
    \end{tabular}

    \textbf{Beak} (7)\\
    Piercing 1\\


    \subsection*{Seal}

    Along cold shores and drifting ice, the seal moves with effortless grace between water and land. In the sea it is swift and elusive, diving and turning with fluid precision; on land it appears slower, yet remains alert and wary. Seals are not natural aggressors, but when cornered or defending their young, they react with sudden ferocity—biting and thrashing with surprising strength. Their presence often signals rich hunting grounds, but also hidden dangers beneath the surface. In darker interpretations, they feel watchful and knowing, eyes lingering just above the waterline before slipping silently out of sight.

    \begin{tabular}{@{}ll@{}}
        \textbf{\trans{wounds_label}:} & 5 \\
        \textbf{\trans{movement_label}:} & 5 \\
        \textbf{\trans{strength_label}:} & 1 \\
        \textbf{\trans{dexterity_label}:} & 2 \\
        \textbf{\trans{mind_label}:} & 2 \\
    \end{tabular}

    \textbf{Bite} (5)\\
    Piercing 0\\


    \subsection*{Rhinoceros}

    A heavily built and short-tempered giant, the rhinoceros embodies raw impact, stubborn resilience, and explosive aggression. Its poor eyesight is offset by acute hearing and scent, reacting instantly to perceived threats. Once it commits, it charges with unstoppable force, goring and crushing anything in its path. A rhinoceros does not maneuver or hesitate—it breaks through. In darker interpretations, it feels like a living battering ram, driven by blind momentum and unyielding instinct.

    \begin{tabular}{@{}ll@{}}
        \textbf{\trans{wounds_label}:} & 10 \\
        \textbf{\trans{movement_label}:} & 4 \\
        \textbf{\trans{strength_label}:} & 3 \\
        \textbf{\trans{dexterity_label}:} & 3 \\
        \textbf{\trans{mind_label}:} & 1 \\
    \end{tabular}

    \textbf{Charge} (10)\\
    Piercing 2\\


    \subsection*{Schlange}

    A silent and elusive predator, the venomous snake embodies patience, precision, and hidden lethality. It relies on camouflage and stillness, striking with lightning speed when prey comes within reach. Its venom does the true killing—weakening, paralyzing, or slowly breaking the victim from within. Snakes avoid open confrontation, preferring ambush and retreat, but a single mistake near them can prove fatal. In darker interpretations, they feel cold and calculating, as if every movement is deliberate and every strike inevitable.

    \begin{tabular}{@{}ll@{}}
        \textbf{\trans{wounds_label}:} & 2 \\
        \textbf{\trans{movement_label}:} & 5 \\
        \textbf{\trans{strength_label}:} & 1 \\
        \textbf{\trans{dexterity_label}:} & 3 \\
        \textbf{\trans{mind_label}:} & 1 \\
    \end{tabular}

    \textbf{Bite} (8)\\
    Poison 1\\


    \subsection*{Rat}

    In the dark and hidden places, the rat thrives where others falter. It slips through cracks, nests in walls, and gathers in numbers that turn nuisance into threat. Alone it is timid and easily driven off, but in groups it becomes bold—gnawing, biting, and overwhelming through sheer persistence. Rats spread filth, consume supplies, and carry unseen dangers with them. In darker interpretations, they feel like a creeping infestation, a living sign of decay that grows stronger the longer it is ignored.

    \begin{tabular}{@{}ll@{}}
        \textbf{\trans{wounds_label}:} & 2 \\
        \textbf{\trans{movement_label}:} & 5 \\
        \textbf{\trans{strength_label}:} & 1 \\
        \textbf{\trans{dexterity_label}:} & 3 \\
        \textbf{\trans{mind_label}:} & 1 \\
    \end{tabular}

    \textbf{Bite} (5)\\
    Bleeding 1\\


    \subsection*{Scorpion}

    Hidden beneath sand, stone, or shadow, the scorpion waits with unnerving stillness. It reacts in an instant when disturbed, its stinger striking with precise, practiced speed. The venom does not need brute force—it weakens, paralyzes, or ends the fight outright. It does not chase; it punishes intrusion. In harsher interpretations, it feels ancient and indifferent, perfectly suited to a world where hesitation means death.

    \begin{tabular}{@{}ll@{}}
        \textbf{\trans{wounds_label}:} & 3 \\
        \textbf{\trans{movement_label}:} & 5 \\
        \textbf{\trans{strength_label}:} & 1 \\
        \textbf{\trans{dexterity_label}:} & 3 \\
        \textbf{\trans{mind_label}:} & 1 \\
    \end{tabular}

    \textbf{Bite} (6)\\
    Poison 1\\
    \textbf{Sting} (8)\\
    Poison 2\\


    \subsection*{Frog}

    Bright colors in the undergrowth signal danger long before the frog itself moves. Small and easily overlooked, it relies on potent toxins rather than strength, turning even a careless touch into a threat. It does not hunt through pursuit but through proximity—its presence alone can be enough to harm. In dense jungles or humid environments, it becomes part of the landscape, a hidden hazard among leaves and water. In darker interpretations, it feels almost unnatural, its vivid appearance a deliberate warning that something beneath the surface is deeply wrong.

    \begin{tabular}{@{}ll@{}}
        \textbf{\trans{wounds_label}:} & 2 \\
        \textbf{\trans{movement_label}:} & 4 \\
        \textbf{\trans{strength_label}:} & 1 \\
        \textbf{\trans{dexterity_label}:} & 1 \\
        \textbf{\trans{mind_label}:} & 1 \\
    \end{tabular}

    \textbf{Tong} (3)\\
    Poison 1\\


    \subsection*{Tiger}

    A solitary and elusive apex predator, the tiger embodies stealth, patience, and explosive violence. Unlike pack hunters, it relies on perfect timing—stalking unseen, closing distance in silence, and ending the hunt in a single decisive strike. Tigers are territorial and unpredictable, vanishing into their surroundings as easily as they emerge. In darker interpretations, they feel almost spectral—silent, inevitable, and terrifyingly precise.

    \begin{tabular}{@{}ll@{}}
        \textbf{\trans{wounds_label}:} & 6 \\
        \textbf{\trans{movement_label}:} & 4 \\
        \textbf{\trans{strength_label}:} & 2 \\
        \textbf{\trans{dexterity_label}:} & 3 \\
        \textbf{\trans{mind_label}:} & 2 \\
    \end{tabular}

    \textbf{Bite} (10)\\
    Piercing 1, Bleeding 1\\
    \textbf{Claws} (10)\\
    Piercing 1\\


    \subsection*{Jellyfish}

    Drifting through the water with quiet grace, the jellyfish appears almost harmless—until contact is made. Its translucent body conceals trailing tentacles armed with potent venom, reacting instantly to touch. It does not pursue or hunt in the traditional sense; instead, it turns its surroundings into a passive hazard, where movement itself becomes dangerous. In numbers, they form near-invisible barriers that are difficult to detect until it is too late. In darker interpretations, they feel like living traps of the sea—silent, indifferent, and unavoidable.

    \begin{tabular}{@{}ll@{}}
        \textbf{\trans{wounds_label}:} & 2 \\
        \textbf{\trans{movement_label}:} & 4 \\
        \textbf{\trans{strength_label}:} & 1 \\
        \textbf{\trans{dexterity_label}:} & 2 \\
        \textbf{\trans{mind_label}:} & 1 \\
    \end{tabular}

    \textbf{Sting} (4)\\
    Poison 1\\


    \subsection*{Ox}

    A powerful and stubborn beast, the bull embodies raw force, endurance, and territorial fury. Whether wild or driven to rage, it reacts with sudden, violent charges, using its horns to gore and toss aside anything in its path. Unlike predators, it does not hunt—it confronts, driving intruders away through sheer momentum and intimidation. Oxen, while more controlled, still possess immense strength and can become dangerous when stressed or overburdened. In darker interpretations, the bull becomes a symbol of blind wrath and unstoppable momentum, a living force that cannot be reasoned with once unleashed.

    \begin{tabular}{@{}ll@{}}
        \textbf{\trans{wounds_label}:} & 8 \\
        \textbf{\trans{movement_label}:} & 4 \\
        \textbf{\trans{strength_label}:} & 2 \\
        \textbf{\trans{dexterity_label}:} & 3 \\
        \textbf{\trans{mind_label}:} & 2 \\
    \end{tabular}

    \textbf{Ram} (10)\\
    Piercing 1\\


    \subsection*{Elephant}

    A towering giant of the natural world, the elephant embodies immense strength, resilience, and unstoppable momentum. Usually calm and deliberate, it becomes a devastating force when threatened or provoked—charging with crushing weight and sweeping aside anything in its path. Its size alone reshapes the battlefield, turning terrain into an advantage. Elephants are not inherently aggressive, but once angered or panicked, they are nearly impossible to stop. In darker interpretations, they feel like living engines of destruction—slow to rouse, but catastrophic when unleashed.

    \begin{tabular}{@{}ll@{}}
        \textbf{\trans{wounds_label}:} & 10 \\
        \textbf{\trans{movement_label}:} & 4 \\
        \textbf{\trans{strength_label}:} & 4 \\
        \textbf{\trans{dexterity_label}:} & 2 \\
        \textbf{\trans{mind_label}:} & 2 \\
    \end{tabular}

    \textbf{Choke} (8)\\
    The elephant chokes the victim with its trunk.\\
    \textbf{Stomp} (16)\\
    The elephant’s full weight bears down on the victim. The attack has Piercing 2, but can only be used if the elephant has already been adjacent to the target for one round.\\


    \subsection*{Eagle}

    High above the land, the eagle circles with commanding vision, observing every movement below. It does not waste effort—waiting for the precise moment before diving in a swift, lethal descent. Its talons strike with accuracy and force, seizing prey before it can react. Even when not attacking, its presence dominates the sky, a constant watcher that cannot be easily approached or escaped. In darker interpretations, it feels like an unblinking eye of the heavens, distant, patient, and utterly decisive when it acts.

    \begin{tabular}{@{}ll@{}}
        \textbf{\trans{wounds_label}:} & 2 \\
        \textbf{\trans{movement_label}:} & 10 \\
        \textbf{\trans{strength_label}:} & 1 \\
        \textbf{\trans{dexterity_label}:} & 4 \\
        \textbf{\trans{mind_label}:} & 2 \\
    \end{tabular}

    \textbf{Claws} (7)\\
    Bleeding 1\\
    \textbf{Beak} (7)\\
    Piercing 1\\


    \subsection*{Bat}

    From the rafters and cave ceilings, bats descend in restless motion, guided more by sound than sight. They move in erratic swarms, disorienting and unnerving, their sudden proximity enough to startle even seasoned travelers. Alone they are fragile, but in numbers they overwhelm—clawing, biting, and filling the air with chaotic movement. Their presence often marks decay, darkness, or forgotten places. In darker interpretations, they feel like living shadows, erupting from the black and vanishing just as quickly.

    \begin{tabular}{@{}ll@{}}
        \textbf{\trans{wounds_label}:} & 2 \\
        \textbf{\trans{movement_label}:} & 5 \\
        \textbf{\trans{strength_label}:} & 1 \\
        \textbf{\trans{dexterity_label}:} & 4 \\
        \textbf{\trans{mind_label}:} & 1 \\
    \end{tabular}

    \textbf{Bite} (2)\\
    Piercing 2\\


    \subsection*{Boar}

    A brutal and fearless force of the wild, the boar embodies aggression, resilience, and territorial dominance. When threatened or cornered, it charges with explosive speed, using its tusks to gore and trample anything in its path. Boars do not retreat easily—they commit fully, fighting through injury and pain with relentless fury. In darker interpretations, they appear almost unstoppable, driven by blind rage and an unyielding instinct to destroy intruders.

    \begin{tabular}{@{}ll@{}}
        \textbf{\trans{wounds_label}:} & 8 \\
        \textbf{\trans{movement_label}:} & 4 \\
        \textbf{\trans{strength_label}:} & 3 \\
        \textbf{\trans{dexterity_label}:} & 1 \\
        \textbf{\trans{mind_label}:} & 2 \\
    \end{tabular}

    \textbf{Ram} (10)\\
    Piercing 1\\


    \subsection*{Camel}

    A hardy beast of harsh lands turned volatile and dangerous, the camel embodies endurance twisted into aggression. Normally resilient and deliberate, it can become fiercely territorial or unpredictable under stress, lashing out with powerful bites, trampling kicks, and sudden charges. Its endurance allows it to outlast opponents, pressing the attack longer than expected. Camels are difficult to intimidate once enraged, their stubborn nature turning into relentless hostility. In darker interpretations, they feel uncanny—enduring, spiteful, and disturbingly persistent, as if the desert itself refuses to yield.

    \begin{tabular}{@{}ll@{}}
        \textbf{\trans{wounds_label}:} & 8 \\
        \textbf{\trans{movement_label}:} & 4 \\
        \textbf{\trans{strength_label}:} & 2 \\
        \textbf{\trans{dexterity_label}:} & 2 \\
        \textbf{\trans{mind_label}:} & 2 \\
    \end{tabular}

    \textbf{Bite} (6)\\
    Piercing 1\\


    \subsection*{Gorilla}

    A massively built and intelligent primate, the gorilla embodies controlled strength, social dominance, and explosive violence. Normally calm within its group, it reacts with terrifying force when threatened, defending territory or kin with overwhelming physical power. Displays of dominance—beating its chest, roaring, charging—are often warnings, but once it commits, the attack is swift and brutal. A gorilla does not hunt like a predator; it confronts and crushes threats directly. In darker interpretations, it feels almost sentient in its fury—calculated, protective, and unstoppable when roused.

    \begin{tabular}{@{}ll@{}}
        \textbf{\trans{wounds_label}:} & 6 \\
        \textbf{\trans{movement_label}:} & 4 \\
        \textbf{\trans{strength_label}:} & 2 \\
        \textbf{\trans{dexterity_label}:} & 3 \\
        \textbf{\trans{mind_label}:} & 3 \\
    \end{tabular}

    \textbf{Grip} (8)\\
    The victim is held fast and can only break free with a Strength test.\\


    \subsection*{Owl}

    In the quiet hours, the owl moves where sound itself seems to fade. Its flight is nearly silent, wings cutting through the dark without warning. Guided by exceptional hearing, it locates movement others would miss, striking with sudden precision from above. It does not linger—appearing only long enough to seize what it seeks before vanishing again into shadow. In darker interpretations, it feels like a watcher of the night, unseen yet ever-present, a whisper of motion just beyond perception.

    \begin{tabular}{@{}ll@{}}
        \textbf{\trans{wounds_label}:} & 2 \\
        \textbf{\trans{movement_label}:} & 8 \\
        \textbf{\trans{strength_label}:} & 1 \\
        \textbf{\trans{dexterity_label}:} & 3 \\
        \textbf{\trans{mind_label}:} & 4 \\
    \end{tabular}

    \textbf{Beak} (6)\\
    Piercing 0\\


    \subsection*{Goose}

    By still water and open fields, the goose holds its ground with unexpected defiance. What seems like a harmless creature quickly turns confrontational when approached, hissing, spreading its wings, and advancing without hesitation. It defends territory and kin with startling boldness, striking with its beak and battering with powerful wings. Retreat is not its first instinct—challenge is. In darker interpretations, it feels oddly relentless, a small but furious guardian that refuses to yield, no matter the odds.

    \begin{tabular}{@{}ll@{}}
        \textbf{\trans{wounds_label}:} & 2 \\
        \textbf{\trans{movement_label}:} & 6 \\
        \textbf{\trans{strength_label}:} & 1 \\
        \textbf{\trans{dexterity_label}:} & 2 \\
        \textbf{\trans{mind_label}:} & 3 \\
    \end{tabular}

    \textbf{Bite} (8)\\
    Bleeding 1\\


    \subsection*{Wild Dog}

    A relentless pack hunter driven by coordination, endurance, and instinct. Wild dogs do not rely on strength alone—they win through teamwork, communication, and exhaustion of their prey. They harry targets over long distances, forcing mistakes and isolating the weak before closing in together. Individually manageable, but in numbers they become a fast-moving, disciplined threat. In darker interpretations, they feel eerily synchronized—like a single will spread across many bodies.

    \begin{tabular}{@{}ll@{}}
        \textbf{\trans{wounds_label}:} & 6 \\
        \textbf{\trans{movement_label}:} & 4 \\
        \textbf{\trans{strength_label}:} & 2 \\
        \textbf{\trans{dexterity_label}:} & 2 \\
        \textbf{\trans{mind_label}:} & 2 \\
    \end{tabular}

    \textbf{Bite} (6)\\
    Bleeding 1\\
    \textbf{Call to the pack} (6)\\
    The wild dog calls to its pack. After 1d3 combat rounds, additional wild dogs arrive equal to the number of successes on this roll.\\


    \subsection*{Giant Crab}

    Across tidal flats and along rocky shores, the giant crab advances with unsettling purpose. Its heavy shell turns aside lesser blows, while massive claws snap with crushing force, capable of breaking bone or armor alike. It does not rush blindly—each movement is measured, closing distance while guarding itself with raised pincers. In its domain, footing is treacherous, and escape is hindered by terrain it knows well. In darker interpretations, it feels like a living bastion of the coast—unyielding, alien, and perfectly adapted to drag the unwary into a slow, inevitable end.

    \begin{tabular}{@{}ll@{}}
        \textbf{\trans{wounds_label}:} & 6 \\
        \textbf{\trans{movement_label}:} & 6 \\
        \textbf{\trans{strength_label}:} & 2 \\
        \textbf{\trans{dexterity_label}:} & 2 \\
        \textbf{\trans{mind_label}:} & 1 \\
    \end{tabular}

    \textbf{Claws} (10)\\
    \textit{Dad-a-chum? Dum-a-chum? Ded-a-chek? Did-a-chick?}\\


    \subsection*{Wolf}

    A predatory hunter of the wild, the wolf embodies patience, instinct, and coordination. Rarely alone, it stalks its prey in packs, testing for weakness before committing to the kill. Wolves avoid unnecessary risk, circling, harassing, and striking when advantage is certain. In harsher interpretations, they may appear unnaturally intelligent, relentless, or disturbingly fearless.

    \begin{tabular}{@{}ll@{}}
        \textbf{\trans{wounds_label}:} & 4 \\
        \textbf{\trans{movement_label}:} & 5 \\
        \textbf{\trans{strength_label}:} & 2 \\
        \textbf{\trans{dexterity_label}:} & 3 \\
        \textbf{\trans{mind_label}:} & 2 \\
    \end{tabular}

    \textbf{Bite} (7)\\
    Bleeding 1\\


    \subsection*{Spider}

    A patient and methodical predator, the venomous spider embodies stillness, precision, and entrapment. It waits unseen, concealed in webs or hidden corners, striking only when prey is caught or vulnerable. Its venom weakens or paralyzes, ensuring the victim cannot escape once seized. Unlike hunters that pursue, the spider controls its domain—turning space itself into a trap. In darker interpretations, it feels calculating and alien, a silent architect of death whose presence is often noticed too late.

    \begin{tabular}{@{}ll@{}}
        \textbf{\trans{wounds_label}:} & 2 \\
        \textbf{\trans{movement_label}:} & 2 \\
        \textbf{\trans{strength_label}:} & 1 \\
        \textbf{\trans{dexterity_label}:} & 3 \\
        \textbf{\trans{mind_label}:} & 1 \\
    \end{tabular}

    \textbf{Bite} (4)\\
    Poison 1, Piercing 1\\


    \subsection*{Giant Octopus}

    From the depths, something vast and patient unfolds. The giant octopus is a master of concealment, blending seamlessly with rock and reef until it chooses to act. Its intelligence is unsettling—testing, probing, learning. When it strikes, it does so with sudden reach, ensnaring prey in powerful limbs and pulling it into a crushing, inescapable grip. The water becomes its weapon, obscuring vision and direction alike. In darker interpretations, it feels ancient and watchful, a mind lurking below, waiting for the right moment to claim what drifts too close.

    \begin{tabular}{@{}ll@{}}
        \textbf{\trans{wounds_label}:} & 10 \\
        \textbf{\trans{movement_label}:} & 4 \\
        \textbf{\trans{strength_label}:} & 3 \\
        \textbf{\trans{dexterity_label}:} & 3 \\
        \textbf{\trans{mind_label}:} & 2 \\
    \end{tabular}

    \textbf{Grip} (10)\\
    The victim is grappled and suffers 2 Wounds each combat round. It can break free on its turn with a Strength test.\\


    \subsection*{Bear}

    A massive and resilient force of nature, the bear embodies raw strength, endurance, and territorial fury. Unlike swift predators, it does not rely on speed or subtlety—once roused, it overwhelms through sheer power and persistence. Bears are often indifferent until provoked, but when threatened or defending their domain, they become relentless and devastating. In darker interpretations, they feel almost unstoppable—an embodiment of primal wrath given form.

    \begin{tabular}{@{}ll@{}}
        \textbf{\trans{wounds_label}:} & 8 \\
        \textbf{\trans{movement_label}:} & 3 \\
        \textbf{\trans{strength_label}:} & 3 \\
        \textbf{\trans{dexterity_label}:} & 2 \\
        \textbf{\trans{mind_label}:} & 2 \\
    \end{tabular}

    \textbf{Bite} (8)\\
    Bleeding 2\\
    \textbf{Claws} (8)\\
    Bleeding 1\\


    \subsection*{Crocodile}

    A patient and ancient ambush predator, the crocodile embodies stillness, precision, and sudden, overwhelming violence. It waits half-submerged, barely visible, striking in an instant when prey comes too close to the water’s edge. Its grip is relentless—once seized, escape is unlikely. Crocodiles do not pursue far; they control the boundary between land and water, turning safe ground into a deadly trap. In darker interpretations, they feel primordial and calculating, as if the water itself has teeth.

    \begin{tabular}{@{}ll@{}}
        \textbf{\trans{wounds_label}:} & 6 \\
        \textbf{\trans{movement_label}:} & 4 \\
        \textbf{\trans{strength_label}:} & 2 \\
        \textbf{\trans{dexterity_label}:} & 2 \\
        \textbf{\trans{mind_label}:} & 2 \\
    \end{tabular}

    \textbf{Bite} (10)\\
    The victim can only break free from the bite with a Strength test. The crocodile will try to drag its victim underwater whenever possible.\\


    \subsection*{Shark}

    Beneath the surface, the shark moves with silent certainty, sensing what cannot be seen. It does not rush blindly—it circles, tests, and closes in with deliberate intent. A sudden strike ends the encounter in a burst of motion and blood, vanishing again into the depths just as quickly. In its domain, escape is uncertain and footing is lost. In darker interpretations, it feels like an unseen force of the deep, drawn to weakness and disturbance, as if the water itself has turned against its prey.

    \begin{tabular}{@{}ll@{}}
        \textbf{\trans{wounds_label}:} & 8 \\
        \textbf{\trans{movement_label}:} & 6 \\
        \textbf{\trans{strength_label}:} & 2 \\
        \textbf{\trans{dexterity_label}:} & 3 \\
        \textbf{\trans{mind_label}:} & 1 \\
    \end{tabular}

    \textbf{Bite} (10)\\
    Piercing 1\\


    \subsection*{Hyena}

    A cunning scavenger and opportunistic hunter, the hyena thrives where others falter. It survives through adaptability, endurance, and ruthless pragmatism, often circling conflict rather than initiating it. Hyenas test weakness relentlessly, harassing, outlasting, and exploiting any sign of vulnerability. Their eerie vocalizations and unsettling demeanor give them a reputation that borders on the unnatural. In darker interpretations, they feel mocking and cruel—creatures that revel in decay and the downfall of others.

    \begin{tabular}{@{}ll@{}}
        \textbf{\trans{wounds_label}:} & 4 \\
        \textbf{\trans{movement_label}:} & 4 \\
        \textbf{\trans{strength_label}:} & 1 \\
        \textbf{\trans{dexterity_label}:} & 3 \\
        \textbf{\trans{mind_label}:} & 2 \\
    \end{tabular}

    \textbf{Bite} (8)\\
    Bleeding 1\\


    \subsection*{Lion}

    A dominant apex predator, the lion embodies strength, authority, and lethal precision. Unlike relentless pack hunters, it relies on bursts of overwhelming force, often striking from concealment or claiming territory through sheer presence. A lion does not chase endlessly—it decides when the hunt begins and ends. In harsher interpretations, it appears regal and terrifying, a symbol of power that tolerates no challenge.

    \begin{tabular}{@{}ll@{}}
        \textbf{\trans{wounds_label}:} & 6 \\
        \textbf{\trans{movement_label}:} & 4 \\
        \textbf{\trans{strength_label}:} & 3 \\
        \textbf{\trans{dexterity_label}:} & 2 \\
        \textbf{\trans{mind_label}:} & 2 \\
    \end{tabular}

    \textbf{Bite} (8)\\
    Piercing 1, Bleeding 1\\
    \textbf{Claws} (6)\\
    Piercing 1, Bleeding 2\\


    \section{Magical}



    \subsection*{Werewolf}

    The werewolf is a creature that is half wolf, half human. At full moon, the normal human transforms into a wolf, which can grow up to two steps tall and is extremely aggressive. Anyone can become a werewolf as soon as they are bitten by an already infected werewolf.

    \begin{tabular}{@{}ll@{}}
        \textbf{\trans{wounds_label}:} & 12 \\
        \textbf{\trans{movement_label}:} & 8 \\
        \textbf{\trans{strength_label}:} & 5 \\
        \textbf{\trans{dexterity_label}:} & 3 \\
        \textbf{\trans{mind_label}:} & 3 \\
    \end{tabular}

    \textbf{Bite} (8)\\
    Piercing 1, Bleeding 1\\
    \textbf{Claws} (8)\\
    Bleeding 1\\


    \subsection*{Compost fairy}

    The Compost Fairy is a strange, magical creature. She once lived in the compost of the witch Mare, in the Middle Ages on earth, not far from the town of Aquisgrani. Mare and she share a special bond, which is also magical.

    \begin{tabular}{@{}ll@{}}
        \textbf{\trans{wounds_label}:} & 2 \\
        \textbf{\trans{movement_label}:} & 8 \\
        \textbf{\trans{strength_label}:} & 2 \\
        \textbf{\trans{dexterity_label}:} & 5 \\
        \textbf{\trans{mind_label}:} & 3 \\
        \textbf{\trans{resistances_label}:} & ['Magic'] \\
    \end{tabular}

    \textbf{Energy Orb} (2)\\
    Piercing 2, Shocked 1\\


    \subsection*{Shimmer whale (leviatan)}

    The Shimmer whale resembles an earthly blue whale in shape, but is slimmer and has longer, almost wing-like side fins. What makes it special is its skin: it is not gray, but has a pearlescent, semi-transparent surface.

Depending on the incidence of light and the magical saturation of the environment, its skin refracts light into all colors of the spectrum. Hence the name. When an iridescent whale breaks the surface, it looks like a living rainbow rising out of the water.
Scholars believe that these whales not only feed on krill, but also absorb the light of the moon and stars when they come to the surface at night.

Shimmer whales avoid shallow coastal waters. They travel through the deep oceans, far away from the routes of merchant ships.

They are most often sighted in the Southern Ocean, in the cold currents far from the heat of the jungle, or in the mystical waters around abandoned island archives.

They travel in small family groups (pods). It is said that their song can calm storms or drive those who listen to it for too long mad.

Since hunting them is extremely dangerous (shimmer whales rarely defend themselves, but are often protected by sea elementals or mermaids) and the animals are rare, the price is enormous.

The Ancatir consider hunting shimmer whales a sacrilege. They only use oil that comes from naturally stranded whales (“gift of the tides”). Alchemists who use “bloody oil” are often expelled from the city in elven enclaves.

The Shimmer Whale is one of the most fascinating and peaceful giants of the seas around Tirakan. It is not just an animal, but a living anomaly closely connected to the magical currents of the oceans.

Shimmer whales never attack first. If killed, all involved suffer a curse of the sea (+2 exhaustion per day at sea).

We saw it at new moon. It glowed beneath the keel like a sunken city. When we threw the harpoons, the beast did not scream. It began to sing. A sound so deep that the wood of my ship splintered and two of my men simply jumped into the water, smiles on their faces. We killed it, yes. But the oil... it burns in the lamps of my cabin, and I swear I see the faces of those who jumped in the shadows.
From the logbook of the whaler 'Haken-Ulf'

    \begin{tabular}{@{}ll@{}}
        \textbf{\trans{wounds_label}:} & 80 \\
        \textbf{\trans{movement_label}:} & 20 \\
        \textbf{\trans{strength_label}:} & 16 \\
        \textbf{\trans{dexterity_label}:} & 12 \\
        \textbf{\trans{mind_label}:} & 11 \\
    \end{tabular}

    \textbf{Sonar pulse} (16)\\
    The whale emits an extremely loud, directional sound pulse.

Targets must pass a Resistance check (WI). If they fail, they are stunned for 1D6 rounds and suffer a -4 penalty to all checks.\\
    \textbf{Powerful fin stroke} (14)\\
    A blow with its massive tail fin that whips the water.

3d6+16 damage. Hits all targets within a 5-meter radius behind the whale.\\
    \textbf{Iridescent glow} (14)\\
    The whale flashes its skin in a dazzling spectrum of colors.

All attackers in close range must pass a Perception check. If they fail, they are blinded for the next round and cannot attack.\\
    \textbf{Deep dive} (12)\\
    The whale retreats into the depths.

Increases the RS by an additional +4 shields for one round while it leaves the combat area.\\

\end{multicols}